\minitopic{从相信到置信度}传统知识论常把相信视为“有或没有”，形式认识论则用\term{置信度}{credence}表示程度。若置信度满足概率公理，则 $0\leq P(A)\leq1$，必然命题概率为1，互斥命题的析取概率等于概率之和。概率并不自动等于客观频率；它在这里首先约束一组信心如何彼此一致。 若同时高度相信 $A$、$B$，却极低相信 $A\land B$，信念组合可能不融贯。荷兰赌论证试图说明，违反概率公理的置信度会使主体接受一组无论结果如何都净亏损的赌注。它把认识规范连接到实际一致性，但也受限于效用、赌注接受和理想化假设\parencite{joyce1998}。

\minitopic{二值信念与程度信念并非竞争替代}“我相信会下雨”与“我对下雨有六成把握”回答不同问题。前者可进入推理、断言和行动承诺，后者表达不确定的强弱；把六成以上一律转成相信，会遇到阈值为何设在此处、多个分别可信命题的合取为何迅速跌破阈值等困难。反过来，只有程度而没有接受状态，也难解释主体为何用某前提继续推理或公开负责。较好的模型把二者连接为任务相关关系：置信度影响何时接受，接受又受行动风险、信息成本与推理用途约束。

\minitopic{一致性只是理性的一层}一个主体可以把每项概率安排得完全一致，却系统忽略反证、选择偏见来源或拒绝更新；另一个主体可能因计算能力有限而略有数值不一致，却敏锐追踪证据。形式一致性识别的是信念组合中的一种缺陷，不等于全部认识理性。评价时应分开内部约束、与证据的响应关系、模型同世界的联系和主体的认知限制。规范理论若要求现实主体精确存储巨大概率分布，还需说明近似、启发式与有限计算怎样在资源约束下仍可合理。

\minitopic{条件概率与更新}条件概率定义为：
\[
P(H\mid E)=\frac{P(H\land E)}{P(E)}.
\]
贝叶斯定理把假设 $H$ 在证据 $E$ 后的概率写为：
\[
P(H\mid E)=\frac{P(E\mid H)P(H)}{P(E\mid H)P(H)+P(E\mid\neg H)P(\neg H)}.
\]
代入开篇数据：
\[
P(H\mid +)=\frac{0.90\times0.01}{0.90\times0.01+0.05\times0.99}\approx0.154.
\]
阳性使患病概率从1\%升至约15.4\%，是重要证据，却远非90\%。90\%是“患病时阳性”的概率，不是“阳性时患病”的概率。

\minitopic{证据强度与似然}证据支持哪个假设，可比较似然 $P(E\mid H)$ 与 $P(E\mid\neg H)$。二者之比越大，$E$越有利于$H$。但后验概率仍依赖先验概率。低先验假设即使得到有利证据，也未必成为最可信解释。 贝叶斯条件化要求新证据被赋予确定值，并把其他概率按规则调整。现实主体常不确定证据本身、忘记旧数据或无法列出全部假设。形式模型是规范性理想和诊断工具，不是对心理过程的逐字描述\parencite{howson2006}。

\minitopic{旧证据与选择性报告}若一个理论在看到数据后才被设计来拟合数据，$P(E\mid H)$ 很高却未必说明预测成功。模型复杂度、独立检验和事先登记有助于防止把噪声当证据。只公布阳性结果会扭曲读者的更新输入；这说明形式理性依赖社会性的信息生产规则。 概率也不直接等于知识。极高概率的彩票命题可能仍不被知道；正常知觉知识则未必伴随精确数值。形式工具最擅长澄清证据比较和置信度变化，不应未经论证取代知识概念。

\begin{argumentbox}
基础率忽视的错误：把 $P(E\mid H)$ 当成 $P(H\mid E)$。纠正步骤是列出一批代表性个体：一万人中约100人患病，90人真阳性；9900名健康者中约495人假阳性。因此阳性者共585人，其中约90人患病。
\end{argumentbox}

\begin{comparison}
墨家“以故生、以理成、以类取”强调理由、规则和类比在论辩中的作用。它不是概率论的早期版本；可比之处仅在于把判断置于可公开检验的推理规范中。现代形式化增加了数量约束，却同样需要回答输入从何而来、分类是否恰当。
\end{comparison}

\minitopic{先验从哪里来}贝叶斯方法并不消灭主观判断，而是把它放在先验概率中公开。先验可来自基础率、既有研究、对称原则或较宽的无信息分布。批评者担心不同先验导致任意结论；支持者回应，足够强且有区分力的证据常使合理先验逐渐汇合。双方都承认，极端先验可能对证据“免疫”。 比较分析应报告先验敏感性：在一组合理先验下，结论是否稳定？若只在特定先验下成立，就应把依赖关系写进结论。透明性不是假装没有判断，而是让判断的位置可见。

\minitopic{条件独立与重复计算}连续收到两条证据时，只有在给定假设后它们条件独立，才可简单相乘似然。两家媒体若引用同一消息源、两个医学检测若受同一生理因素影响，就不能当成独立确认。忽视相关性会产生过度自信。 自然频率有助于检查：把百分比转换为一千或一万人中的人数，画出真阳性、假阳性、真阴性和假阴性。形式工具的教育价值不是炫示公式，而是迫使推理者保留分母、方向和依赖结构。

\minitopic{案例工作坊：两次检测}假定第二次检测与第一次条件独立且参数相同。第一次阳性后的患病概率约15.4\%。第二次也阳性时，可用第一次后验作为新先验，得到约76.7\%。若两次使用同一污染样本，条件独立假设失败，这个数值会高估支持。答题时必须先写假设，再计算结果。 形式答案还应转回行动语言。76.7\%是否足以治疗，取决于治疗副作用、进一步检测成本和延误风险。概率描述信念与证据，决策还需要效用；知识判断则可能关心个案联系与来源，而不只数值门槛。
